\section{项目应用情况}
\subsection{落地应用}

项目当前处于方案设计阶段，具体竞赛演示场景尚未最终确定，尚无已落地应用或企业试用记录。首期应用定位为面向嵌入式推理开发的FPGA原型：使用现有试验箱，验证模型从电脑端准备、转换，到板上执行与结果回传的完整过程。应用扩展按“基础推理、静态视觉、感知接入”的顺序推进。\cite{technicalreport}

首个可行性算例为\(\ModelImageHeight\times\ModelImageWidth\)灰度数字图像，采用\ModelInputSize{}$\rightarrow$\ModelHiddenSize{}$\rightarrow$\ModelOutputSize{}的小型\ModelFamily{}。电脑端准备来源清晰的训练权重、固定预处理、校准样本、独立测试集及逐层参考结果，板端仅执行推理。最终类别由\ModelOutputOperation{}确定；未经过概率变换的分数按分类分数展示。该算例用于检验计算系统的连贯性，不作为最终行业应用的定稿。

{\small\renewcommand{\arraystretch}{\CompactTableStretch}
\noindent\begin{tabularx}{\linewidth}{@{}P{.19\linewidth}Y Y@{}}
\toprule
应用阶段 & 拟验证内容 & 进入下一阶段的依据 \\
\midrule
基础推理 & 固定张量、量化模型与三后端执行 & 逐层结果可对照，完整推理可重复 \\
静态视觉 & 输入图片回放、预处理与类别回传 & 样本和输出可追溯，模型效果独立评估 \\
感知扩展 & 摄像头、传感器与通信适配 & 明确设备、协议和资源条件，再完成接口测试 \\
\bottomrule
\end{tabularx}}

原型的预期价值是为开发者提供从算子描述到执行验证的可检查路径，并展示不同计算后端的适用范围。实际性能、开发工作量和部署成本将通过工程验证记录，形成后续评价依据。题面中的移动数据中心尚未指定平台和接口，方案保留结果输出适配层，具体接入随需求确认后展开。

\subsection{人机交互流程}

首期使用电脑端准备工具和串口交互完成操作。一次推理按以下顺序进行：

\begin{enumerate}
\item 电脑端选择模型与样本，检查形状、量化参数及后端支持情况，生成执行计划、权重和版本标识。
\item 操作者选择\RealTimeSystem{}或\LinuxDistribution{}启动镜像，重启进入相应环境，查看串口返回的初始化和就绪状态。
\item 通过串口加载程序或测试数据，提交带样本标识的任务，由运行时按计划执行并等待各步骤完成。
\item 回传样本标识、分类结果与诊断记录，在电脑端对照参考结果；参数不支持、初始化失败或执行异常时保留错误位置与状态。
\end{enumerate}

后续静态图像与实时采集均通过输入适配层接入，保持运行时接收张量、输出带标识结果的边界。实时摄像头和传感器驱动分别进行现场采集验证。
